\chapter{Il pipeline di sviluppo farmaceutico}
%\begin{figure}[htbp]
    %\centering
    %\includegraphics[width=0.85\textwidth]{figures/pipeline.png}
    %TO-DO
   % \caption{Il pipeline di sviluppo farmaceutico.}
    %\label{fig:pipeline}
%\end{figure}
\section{Le fasi del Drug Development}

Lo sviluppo di un nuovo farmaco è un processo complesso e delicato, 
che si articola in più fasi sequenziali, ciascuna delle quali introduce 
vincoli tecnici, economici e regolatori che rendono l'intero percorso 
lungo, dispendioso e ad alto rischio di fallimento \cite{zhang2025}.

Negli ultimi due decenni si è distinta come principale strategia di discovery la \textbf{Target Based Drug Discovery} \cite{keshTargetBasedScreeningLead2023} (\textbf{TBDD}), processo che ha come inizio l'\textbf{identificazione del target}, 
ovvero la selezione di molecole biologiche, tipicamente  
proteine o acidi nucleici, la cui modulazione risulti rilevante 
per la malattia di interesse (ovvero target \textit{druggable}, trattabili farmalogicamente).
Si tratta del primo filtro a cui vengono sottoposte le molecole candidate ed è proprio il primo "collo di bottiglia" del processo di ricerca. 
Gli svariati approcci di screening utilizzati in questa fase di \textit{identification} possono essere classificati in quattro gruppi principali \cite{rasulTargetIdentificationApproaches2022}:
\begin{itemize}
    \item Approcci "diretti" o basati sull'\textbf{affinità} tra proteine e piccole molecole.
    \item Approcci basati sul \textbf{fenotipo}, usati per effettuare comparazioni tra i profili biologici delle molecole e quelli di farmaci già noti, di cui si conosce il comportamento.
    \item Approcci basati sull'individuazione dei \textbf{geni} specifici responsabili di farmacoresistenza da parte delle molecole.
    \item Approcci \textbf{computazionali}, in grado di fare previsioni sui target sulla base della complementarietà chimica dei candidati.
\end{itemize}
Sebbene molti di questi metodi, come l'\textit{affinity pull-down} e il \textit{whole-genome knockdow} siano ampiamente usati e rimangano i pilastri per la validazione dei target di interesse, questi presentano limiti intrinseci in termini di efficacia e costi. 
%----

Seguono le fasi di \textbf{drug discovery} vero e proprio, che si pongono l'obiettivo di identificare gli \textbf{\textit{hits}}, farmaci in grado di interagire con i target precedentemente individuati.
I geni candidati vengono espressi in grandi quantità in organismi poco complessi tramite la tecnologia del DNA ricombinante, permettendo lo \textbf{High-Throughput Screening} (HTS), ovvero uno screening ad alta velocità che agisce su vaste e variegate librerie chimiche \cite{keshTargetBasedScreeningLead2023}.

Sebbene l'HTS sia più performante in termini di costi rispetto ai metodi manuali, non si è dimostrato in grado di garantire la sicurezza o l'efficacia della molecola su organismi più complessi. 
Per superare tali limiti, emerge la necessità di tecniche di \textbf{Virtual Screening} (VS), che permettono di pre-selezionare \textit{in silico} i composti più promettenti, riducendo drasticamente il numero di test fisici necessari.

\subsection{Screening dei composti}
Il processo di screening rappresenta lo strumento per "setacciare" le librerie chimiche nella ricerca delle giuste interazioni biochimiche. Il \textbf{VS} si distingue dall'HTS per la sua capacità di andare oltre i limiti della ricerca fisica attuabile in laboratorio, valutando computazionalmente composti appartenenti a librerie fisiche o sintetizzabili su richiesta \cite{carlssonStructurebasedVirtualScreening2024}.
Questo approccio si divide in due macro-famiglie:
\begin{itemize}
    \item I metodi \textbf{ligand-based}, che si basano sull'ipotesi che molecole con proprietà simili a ligandi noti presentino attività biologica analoga. Tra questi, la \textbf{pharmacophore searching} identifica composti compatibili al target per caratteristiche chimico-spaziali (come la presenza di donatori e accettori di idrogeno o di gruppi aromatici in determinate posizioni) \cite{mcinnesVirtualScreeningStrategies2007a}.
    \item I metodi \textbf{structure-based}, che sfruttano la struttura tridimensionale del target tramite strumenti come il \textbf{molecular docking} o l'\textbf{ensemble docking}, quest'ultimo è un'evoluzione del primo che valuta il ligando su molteplici conformazioni del recettore, ottenendo così risultati più robusti rispetto all'approssimazione del recettore singolo e rigido.
\end{itemize}
Il virtual screening accelerato dal Machine Learning rappresenta una risposta al crescente e sempre più insostenibile peso computazionale del docking. Questo approccio prevede che un modello di ML venga addestrato sui risultati del docking su un sottoinsieme della libreria (si usano librerie dette "ultra-large") e che impari a predire i punteggi del docking per i composti incognita, senza eseguire il calcolo fisico. Questo approccio può portare a una riduzione di 100 volte il numero di composti da sottoporre a docking esplicito, mantenendo comunque un arricchimento dei composti top-scoring paragonabile allo screening esaustivo \cite{carlssonStructurebasedVirtualScreening2024}.


La \textit{hit discovery} è tuttavia solo il primo passaggio del processo di \textit{lead generation}, o \textit{Hit to Lead} (H2L), che si pone di valutare i composti individuati dalle \textit{hits} non solo per la loro presenza di attività, ma per la qualità della loro interazione con il target. 
In questa fase vengono analizzate potenza, selettività e proprietà fisico-chimiche di base, causando una scrematura che può apportare un miglioramento del legame chimico hit-target dell'ordine di grandezza nanomolare \cite{barcelosLeadOptimizationDrug2022}.
La fase finale di drug discovery è la \textit{Lead Optimization} (LO), che interviene per accentuare caratteristiche desiderate nei composti lead selezionati, o per correggere i difetti strutturali che impedirebbero ad essi di superare i trial clinici.
La LO moderna è un processo razionale guidato da metodi computazionali che includono studi farmacoforici, docking molecolare (metodo utilizzato per prevedere l'orientamento di un ligando su un recettore per la formazione di un complesso stabile), dinamica molecolare (simulazioni con lo scopo di ottenere dati sulle proprietà fisico-strutturali delle molecole) e QSAR. 
In paricolare, il \textbf{QSAR} (Quantitative Structure-Activity Relationship) rappresenta un approccio di modellazione statistica volto a correlare le caratteristiche strutturali con l'attività biologica dei composti, evolvendosi oggi in modelli di Machine Learning capaci di interpretare e prevedere queste caratteristiche per generare risultati nuovi \cite{barcelosLeadOptimizationDrug2022}.

Una volta ottenute le lead ottimizzate si può passare agli \textbf{studi pre-clinici}, condotti su modelli animali per valutare 
sicurezza ed efficacia dei farmaci, prima che questi vengano testati sull'essere umano. 
La fase successiva, quella degli effettivi \textbf{trial clinici}, si 
articola a sua volta in tre stadi progressivi: 
\begin{itemize}
    \item La fase~I verifica la 
sicurezza del composto su un piccolo gruppo di soggetti sani;
    \item La fase~II ne valuta l'efficacia su pazienti affetti dalla patologia interessata;
    \item La fase~III, condotta su popolazioni più ampie, 
costituisce il test definitivo richiesto per l'approvazione regolatoria. 
\end{itemize}  
Al termine di questo percorso, il farmaco approvato entra nella fase 
di \textbf{post-market surveillance}, in cui vengono monitorati nel 
lungo periodo gli effetti collaterali e la stabilità della formulazione.

\section{Attrition Rate e impatto economico}
Il percorso che racchiude queste fasi è caratterizzato da un problema fondamentale: nonostante l'accelerazione in ambito tecnologico degli ultimi anni e la comparsa di nuovi strumenti computazionali ad ausilio del processo, il successo delle fasi cliniche rimane un obiettivo estremamente difficile da raggiungere. Questo fenomeno è descritto dall'\textbf{attrition rate}, ovvero il tasso di riduzione ("logoramento") del numero di candidati farmaci durante l'avanzamento della pipeline.
Delle migliaia di molecole potenziali individuate dalla prima fase di discovery, solo una piccola frazione sopravvive ai test clinici e meno del 10\% dei composti che entrano nel primo step di trial raggiunge l'approvazione finale dalle autorità regolatorie.

Questo elevato tasso di insuccesso si traduce in elevatissimi costi economici.
Lo sviluppo di un nuovo farmaco richiede oggi un investimento medio di circa 
\textbf{2,6 miliardi di dollari americani} e un arco temporale compreso tra 
i dodici e i quindici anni \cite{zhang2025}. 
Inoltre, anche il ritorno economico per le industrie farmaceutiche 
è drasticamente calato: nell'arco di un decennio, per ogni dollaro 
investito in ricerca il ritorno economico è crollato drasticamente, passando da 10 a meno di 2 centesimi 
\cite{freedmanHuntingNewDrugs}.

Le cause di questo elevato tasso di fallimento sono molteplici. 
In primo luogo, molti farmaci per malattie comuni e con meccanismi biologici noti sono già stati scoperti; ciò che rimane da 
esplorare riguarda patologie più rare e complesse, dove l'identificazione 
di molecole target risulta più difficile \cite{freedmanHuntingNewDrugs}. 
Inoltre, la selezione delle molecole nelle fasi di \textit{discovery} può essere basata su parametri toppo semplici e poco esaustivi. Infatti, come evidenziato da \textbf{Kenakin} \cite{kenakinKnowYourMolecule2024}, la mancanza di efficacia clinica e i problemi di sicurezza legati alla tossicità rimangono le cause primarie del logoramento dei candidati farmaci, ed entrambi sono imputabili a una descrizione farmacologica superficiale e poco approfondita dei \textbf{lead} stessi.

\section{Il ruolo dell'AI nella pipeline}
L'integrazione dell'Intelligenza Artificiale (AI) e del Deep Learning (DL)
ha trasformato la pipeline di ricerca farmaceutica da un processo basato su "trial and error" a una percorso più disciplinato, guidato dalla raccolta e dall'analisi dei dati biologici \cite{zhang2025}.
Questi strumenti vengono impiegati trasversalmente, in ciascuna delle diverse fasi più delicate e "logoranti" del processo.

Nella fase di \textbf{target identification}, l'AI sta affiancando 
e in alcuni casi sostituendo i metodi tradizionali di screening, 
prima basati sull'analisi dettagliata dell'affinità tra molecole. 
Il nuovo approccio consiste nel rappresentare i geni per la loro \textit{funzione biologica}, analogamente alle 
tecniche di \textit{word embedding} sviluppate nel Natural Language Processing (NLP). 
Negli step successivi dello sviluppo del composto, il contributo più significativo si osserva nella predizione dei comportamenti di interazione farmaco-proteina (DPI).
Queste stime sono solitamente compiute su uno spazio chimico di scala molto ampia, e richiedono l'analisi di molte interazioni ancora sconosciute. Proprio per questo, i modelli di Machine Learning vengono addestrati in modalità semi-supervisionato per etichettare queste interazioni incognite, integrando efficacemente informazioni eterogenee, quali la struttura chimica del composto, i dati genomici e le reti di interazione.
L'importanza di affinare tali modelli risiede nella necessità di colmare il \textit{knowledge gap} relativo agli effetti collaterali:
è stato infatti stimato che circa 1.000 proteine umane siano potenzialmente responsabili di tossicità \textit{off-target}. 

Identificare precocemente queste interazioni impreviste risulta dunque fondamentale per prevedere e mitigare interazioni nocive già nelle fasi preliminari dello sviluppo in silico \cite{dara2022}.

